\noindent\textbf{Elección del procedimiento.} La pendiente del 6 por ciento significa $\tan\varphi=0{,}06$, por tanto $\sin\varphi=0.06/\sqrt{1+0{,}06^2}$. En descenso terminal sin rodadura, $Mg\sin\varphi=(\rho V^2/2)C_DA$.

\begin{enumerate}

\item Proyectar el peso con el seno exacto del ángulo de la pendiente.

\[\fitmath{\mathrm{Fslope}=\frac{M\cdot g\cdot \mathrm{slope}}{\sqrt{1+{\mathrm{slope}}^{2}}}}\]

\[\fitmath{\mathrm{Fslope}=\frac{(90)\cdot (9{,}81)\cdot (0{,}06)}{\sqrt{1+{(0{,}06)}^{2}}}\approx 52{,}8789\,\mathrm{N}}\]

\item Aislar el producto $C_DA$, no un coeficiente sin área.

\[\fitmath{\mathrm{CDA}=\frac{2\cdot \mathrm{Fslope}}{\rho\cdot {V}^{2}}}\]

\[\fitmath{\mathrm{CDA}=\frac{2\cdot (52{,}87890348)}{(1{,}2)\cdot {(15)}^{2}}\approx 0{,}3916956\,\mathrm{m^{2}}}\]

\end{enumerate}
